\documentclass[11pt]{article}
\usepackage[margin=1in]{geometry}
\usepackage{amsmath,amssymb,amsthm}
\usepackage[hidelinks]{hyperref}

\newtheorem{claim}{Claim}
\newtheorem{remark}{Remark}
\newcommand{\Z}{\mathbb{Z}}
\newcommand{\Graver}{\operatorname{Graver}}

\title{Disproof of TLMC Conjecture 00000000544:\\ the Equality-on-Forests Clause Fails on $K_2$}
\author{TLMC submission --- conjecture 00000000544}
\date{October 2, 2026}

\begin{document}
\maketitle

\section{The conjecture}

Let $G$ be a graph without isolated vertices, $A_G$ its incidence matrix, and
$I_{A_G} \subseteq k[x_1,\dots,x_{|E|}]$ the toric ideal generated by $A_G$, i.e.\ the
kernel of the monomial map
\[
\varphi : k[x_1,\dots,x_{|E|}] \longrightarrow k[t_1^{\pm1},\dots,t_{|V|^{\pm1}}],
\qquad x_e \longmapsto t^{a_e},
\]
where $a_e$ is the column of $A_G$ indexed by the edge $e$. The Graver basis
$\Graver(I_{A_G})$ is the set of all irreducible elements of the lattice
$L_A = \ker_\Z(A) \subseteq \Z^{|E|}$.

Conjecture \textbf{00000000544} asserts that for all graphs without isolated vertices
\begin{equation}\label{eq:conj}
\bigl|\Graver(I_{A_G})\bigr| \;\le\; \sum_{H \subseteq G,\ H \text{ forest}} 2^{|V(H)|},
\qquad\text{with equality attained on forests.}
\end{equation}

\section{The counterexample}

\begin{claim}\label{clm:main}
$G = K_2$ (one edge, two vertices) satisfies the hypothesis but violates the equality
clause of \eqref{eq:conj}. The conjecture is false.
\end{claim}

$K_2$ is a forest and has no isolated vertices, so \eqref{eq:conj} must hold with
equality for $G = K_2$. We compute both sides.

\subsection{Left-hand side: $|\Graver(I_{A_{K_2}})| = 0$}

The (unoriented) incidence matrix of $K_2$ is the $2 \times 1$ matrix
\[
A_{K_2} \;=\; \begin{pmatrix} 1 \\ 1 \end{pmatrix}.
\]
The monomial map is
\[
\varphi : k[x_1] \to k[t_1, t_2], \qquad x_1 \mapsto t_1 t_2 ,
\]
which is injective: $t_1t_2$ is a nonconstant monomial, so
$\ker \varphi = 0$ and $I_{A_{K_2}} = 0$. (Under the oriented convention
$a_1 = (1,-1)^{\mathsf T}$ one gets $x_1 \mapsto t_1t_2^{-1}$, again injective; nothing
below changes.)

Equivalently at the lattice level,
\[
L_A \;=\; \ker_\Z(A_{K_2}) \;=\; \{\, u \in \Z^1 : A_{K_2} u = 0 \,\} \;=\; \{\, u \in \Z : u = 0 \,\} \;=\; \{0\},
\]
since $A_{K_2}u = (u, u)^{\mathsf T}$. The lattice $\{0\}$ has no nonzero elements and
therefore no irreducible ones:
\[
\Graver(I_{A_{K_2}}) = \varnothing, \qquad \bigl|\Graver(I_{A_{K_2}})\bigr| = 0 .
\]

\subsection{Right-hand side: $\mathrm{RHS} \ge 4 > 0$}

The graph $H = G = K_2$ is itself a forest subgraph of $G$ with $|V(H)| = 2$, so its
summand alone contributes
\[
2^{|V(H)|} = 2^2 = 4 .
\]
Hence the right-hand side of \eqref{eq:conj} is at least $4$ under every reading of
``$H \subseteq G$''. For instance, with $H$ ranging over spanning subgraphs there are
exactly two of them, the edgeless graph and $K_2$ itself, both forests on $2$ vertices,
giving $\mathrm{RHS} = 2^2 + 2^2 = 8$.

\subsection{Conclusion}

\[
0 \;=\; \bigl|\Graver(I_{A_{K_2}})\bigr| \;\neq\; \mathrm{RHS} \;\ge\; 4 ,
\]
so \eqref{eq:conj} fails with equality on the forest $K_2$. Conjecture 00000000544 is
\textbf{false}.

\begin{remark}[robustness]\label{rem:robust}
The failure is not an artefact of the smallest example. For \emph{every} nonempty
forest $G$, the columns of the unoriented incidence matrix $A_G$ are linearly
independent over $\mathbb{Q}$, hence $\ker_\Z(A_G) = \{0\}$, $I_{A_G} = 0$ and
$|\Graver(I_{A_G})| = 0$, while the right-hand side is $\ge 2^2 = 4$ as soon as $G$
has an edge. Thus equality fails on \emph{every} forest with at least one edge;
$K_2$ is the minimal counterexample, and the inequality half of \eqref{eq:conj}
survives here only trivially as $0 \le \mathrm{RHS}$.
\end{remark}

\begin{remark}[boundary]
Single-vertex graphs are excluded by the no-isolated-vertices hypothesis, so $K_2$ is
the smallest admissible graph; there is no smaller boundary case. The verification is
machine-checked: a standalone script (\texttt{reproduce.py}) recomputes the lattice,
the Graver basis and the subgraph sum, and a Lean~4 formalization (\texttt{lean4/},
core only, zero axioms, zero \texttt{sorry}) certifies all numerical assertions.
\end{remark}

\end{document}
